% Generated draft: requires mathematical and source review.
\documentclass[11pt]{article}
\usepackage[margin=25mm]{geometry}
\usepackage{fontspec}
\setmainfont{Latin Modern Roman}
\usepackage{amsmath,amssymb,mathtools,unicode-math}
\setmathfont{Latin Modern Math}
\usepackage{enumitem,needspace,xurl,hyperref}
\setlength{\parindent}{0pt}
\setlength{\parskip}{5pt}
\setlength{\emergencystretch}{3em}
\newcommand{\N}{\mathbb{N}}
\newcommand{\Z}{\mathbb{Z}}
\newcommand{\Q}{\mathbb{Q}}
\newcommand{\R}{\mathbb{R}}
\newcommand{\C}{\mathbb{C}}
\newcommand{\F}{\mathbb{F}}
\newcommand{\A}{\mathbb{A}}
\newcommand{\PP}{\mathbb{P}}
\newcommand{\E}{\mathbb{E}}
\newcommand{\eps}{\varepsilon}
\newcommand{\dd}{\,\mathrm{d}}
\begin{document}
\section*{UnsolvedMath ID: 30003496}
\textbf{Source problem code:} catalogue entry
\subsection{Definitions and mathematical statement}
The supplied record asks for the polynomial identities of a given noncommutative algebra \(A\), particularly a matrix algebra, and for the symbolic calculus resulting from those identities. It does not specify the coefficient domain, associativity or unit conventions, the algebra \(A\), or what constitutes the requested symbolic calculus. Consequently, it does not determine a unique mathematical problem.

A conditional formulation of the polynomial-identity part is as follows. Suppose that \(K\) is a field and \(A\) is an associative \(K\)-algebra. These are additional conventions for this formulation, not hypotheses established by the source. For each integer \(m\geq 1\), let \(K\langle x_1,\ldots,x_m\rangle_+\) denote the noncommutative polynomials over \(K\) with zero constant term. Define
\[
\operatorname{Id}_{K,m}(A)=\left\{f\in K\langle x_1,\ldots,x_m\rangle_+ : f(a_1,\ldots,a_m)=0\ \text{for every }(a_1,\ldots,a_m)\in A^m\right\}.
\]
Here evaluation replaces each variable by its assigned element of \(A\), preserving the order of multiplication in every monomial. The zero-constant-term convention permits evaluation without requiring a multiplicative identity; the record does not prescribe this convention.

\textbf{Requested task, conditional on these conventions.} For the given algebra \(A\), characterize \(\operatorname{Id}_{K,m}(A)\) for every \(m\geq 1\). The matrix-algebra specialization would take \(A=M_n(K)\), the algebra of \(n\times n\) matrices over \(K\), for a specified integer \(n\geq 1\), restricting to noncommutative instances when required. Neither \(K\) nor \(n\) is specified in the record. Construct the symbolic calculus intended by the source from these identities; however, its objects, operations, equivalence relation, and required properties are unspecified, so this second request cannot yet be stated precisely. The conditional formulation above does not resolve that missing target.
\subsection{Short English statement}
The record asks which noncommutative polynomial expressions vanish under every substitution into a given algebra, especially a matrix algebra, and what symbolic calculus follows. A conventional associative formulation is given conditionally, but the coefficient domain, particular algebra, and meaning of symbolic calculus remain unspecified.
\subsection{Sources}
\begin{itemize}
\item Supplied local UnsolvedMath registry record: permanent ID \(30003496\), source problem code \texttt{catalogue entry}, title \emph{Polynomial Identities and Symbolic Calculus for Noncommutative Algebras}; collection: Oberwolfach Reports Open Problems. No individual proposer is supplied; the recorded proposed year is 2017.
\item Supplied report citation: \emph{Algebraic Groups (2018)}, Oberwolfach Reports, pp.\ 1281--1347, DOI \texttt{10.4171/owr/2017/21}; \url{https://doi.org/10.4171/owr/2017/21}. The background describes these reports as documenting mathematical workshops and their open-problem discussions. This is the supplied report-level attribution; no precise passage locating this fragment is supplied.
\item Supplied literature source, labelled ``Oberwolfach Reports 2018'': \url{https://doi.org/10.4171/owr/2017/21}. The literature assessment dated \textbf{2026-08-21} labels the record open while stating that the extracted fragment is underspecified and does not identify a definite claim whose truth has been settled. This is supplied dated curation, not an independently verified literature assessment or source visit; time-sensitive status and cited claims require independent verification.
\end{itemize}
\end{document}
